\section{Using the Vocabulary}
\label{sec:uses}

\subsection{Locating a failure}
When an agent fails, the primitives give a short list of places to look, in order. The order matters: each question rules out a place the next question assumes.

\begin{enumerate}
  \item \textbf{\prim{Context}: was the model shown what it needed?} Read the request that was sent. If the file, the result or the instruction is not in it, the model could not have used it, and nothing downstream is at fault.
  \item \textbf{\prim{Output}: did what it asked for happen?} Compare the action the model requested with what ran. A parser that dropped a call, a sandbox that refused it, a tool that returned a truncated result, and an error message the model could not read are all \prim{Output}.
  \item \textbf{\prim{Control flow}: did the loop stop early or run too long?} A step limit reached before the task was done, a hand-back on a response that still had work in it, and a loop that repeated the same action are \prim{Control flow}.
  \item \textbf{\prim{Model interface}: was the response read correctly?} A tool call emitted in the right format but not recognized, a stream cut off and treated as complete, and a backup model substituted without notice are \prim{Model interface}.
  \item \textbf{\prim{Input}: did the right thing come in?} A result delivered to the wrong turn, an interrupt lost, a person's answer consumed by the harness when the model needed it.
  \item \textbf{Only then, the model.} If the model was shown what it needed, asked for the wrong thing, and its request was carried out as asked, the failure is the model's: its knowledge, its habits, or its size.
\end{enumerate}

Two examples from our own transcripts. A small model, after instruction-tuning, sometimes emitted a second tool call as bare JSON without its \texttt{<tool\_call>} tags. The harness read the JSON as words, found no action, and handed back. Question 4 locates this in \prim{Model interface}, whose parser defines what counts as a call; the fix was on the model side (training on the format) but the diagnosis was the harness's. In the same experiments, four kinds of task could never pass during reinforcement learning because the check compared the folder before and after, and the harness's own memory directory changed the folder on every run. No question in the list locates this, because the fault was in the evaluation, a concern, not in a primitive. Building the harness around a capable model first, until it passed every task in the evaluation, is what let every later failure be attributed to the model swapped in.

\subsection{Reporting a harness}
If $\textit{Agent} = \textit{Harness}(\textit{Model})$, a benchmark score belongs to the pair. Holding GPT-4 Turbo fixed, \citet{yang2024sweagent} report 1.31\% of SWE-bench \citep{jimenez2024swebench} resolved with retrieval-augmented generation and 12.47\% with their agent, and on SWE-bench Lite 11.00\% for a shell-only agent against 18.00\% for SWE-agent, with only the harness changing. Anthropic notes that performance ``can vary significantly based on this scaffolding, even when using the same underlying AI model'' \citep{anthropic2025swebench}; METR found that tailoring a scaffold to one model more than doubled the length of tasks it could complete \citep{metr2025update}, and later that two mature harnesses did not outperform their default scaffolds \citep{metr2026harness}, which suggests the effect is largest where the design decisions differ most. \citet{zhang2026stop} argue that the harness should be disclosed whenever agents are compared, and \citet{kapoor2025hal} evaluate models, scaffolds and benchmarks as separate dimensions. A disclosure needs a schema. We suggest the five primitives: the input sources; what the context holds and how it is fitted to the window; the provider, decoding settings and retry policy; the tools, their output formats and the execution environment; and the loop's shape and stopping rules. The summary table of \S\ref{sec:deconstruction} is such a disclosure for six harnesses.

\subsection{Porting}
Because only \prim{Model interface} knows which model is inside, a harness can be moved to a new model or provider by changing that primitive alone, translating requests and responses to the new format. Our own harness was built around a hosted model and later run with a local one by replacing one function (Appendix~\ref{app:minimal}). A harness loses this property when model-specific details leak into other primitives: a provider's tool-call format assumed in \prim{Output}, or a provider's error strings matched in \prim{Context} management. The deconstruction records such leaks where we found them.

\subsection{Training a model for its harness}
The primitives also say what a model trained for one harness must learn: the context layout it will be shown (\prim{Context}), the format in which it must request actions and the shape of the results it will read (\prim{Output}), and the loop it runs in (\prim{Control flow}). A model post-trained for many harnesses must guess which one it is in; a model trained inside one does not.\footnote{In the companion book, a 0.6-billion-parameter model instruction-tuned on sessions recorded in its own harness rose from 2 to 18 of 32 held-out tasks in that harness, where the same base model post-trained by its developers for general use scored 14. The experiment, its data and its checks are described there \citep{course}; this paper makes no claim that rests on it.}
